% !TEX root = ../matan.tex

\section{Признаки сходимости числовых рядов}

\subsection{Ряды с неотрицательными членами}

Пусть \(a_k\ge 0\). Тогда частичные суммы \(S_n=\sum_{k=1}^{n}a_k\) не убывают. Поэтому ряд \(\sum a_k\) сходится тогда и только тогда, когда последовательность \((S_n)\) ограничена сверху.

\begin{Theorem}{Интегральный признак сравнения}{thm:integral-comparison}
	Пусть \(a_k\ge 0\), а функция \(f:[1,+\infty)\to\mathbb R_+\) монотонно убывает.
	
	\begin{enumerate}
		\item Если начиная с некоторого номера \(a_k\le f(k)\) и несобственный интеграл \(\int_1^{+\infty} f(x)\,dx\) сходится, то ряд \(\sum_{k=1}^{\infty} a_k\) сходится.
		
		\item Если начиная с некоторого номера \(a_k\ge f(k)\) и интеграл \(\int_1^{+\infty} f(x)\,dx\) расходится, то ряд \(\sum_{k=1}^{\infty} a_k\) расходится.
	\end{enumerate}
\end{Theorem}

\begin{proof}
Так как \(f\) убывает, то для \(x\in[k,k+1]\) выполнено
\[
f(k+1)\le f(x)\le f(k).
\]
Отсюда получаем две оценки:
\[
f(k+1)\le \int_k^{k+1} f(x)\,dx\le f(k).
\]

Пусть \(a_k\le f(k)\) при \(k\ge k_0\). Тогда для \(k\ge k_0\)
\[
a_k\le f(k)\le \int_{k-1}^{k} f(x)\,dx.
\]
Следовательно, для \(N\ge k_0\):
\[
\sum_{k=k_0}^{N} a_k
\le
\int_{k_0-1}^{N} f(x)\,dx
\le
\int_1^{+\infty} f(x)\,dx.
\]
Частичные суммы ряда с неотрицательными членами ограничены, значит ряд сходится.

Пусть теперь \(a_k\ge f(k)\) при \(k\ge k_0\). Тогда
\[
a_k\ge f(k)\ge \int_k^{k+1}f(x)\,dx.
\]
Поэтому
\[
\sum_{k=k_0}^{N}a_k
\ge
\int_{k_0}^{N+1}f(x)\,dx.
\]
Если интеграл расходится, правая часть стремится к \(+\infty\), значит частичные суммы ряда тоже неограничены. Следовательно, ряд расходится. 
\end{proof}

\begin{Remark}{О конечном числе первых членов}{rem:first-terms}
	Сходимость или расходимость ряда не меняется, если удалить или добавить конечное число членов. Поэтому в признаках обычно достаточно проверять неравенства начиная с некоторого номера.
\end{Remark}

\begin{Example}{}{ex:integral-test-examples}
	\begin{enumerate}
		\item Ряд \(\sum_{k=1}^{\infty} e^{-k}\) сходится, потому что сходится интеграл \(\int_1^{+\infty} e^{-x}\,dx\).
		
		\item Для любого \(m>0\) ряд \(\sum_{k=1}^{\infty} k^m e^{-k}\) сходится. Действительно, функция \(x^m e^{-x}\) начиная с некоторого места убывает, а интеграл \(\int_1^{+\infty} x^m e^{-x}\,dx\) сходится.
		
		\item Для степенного ряда \(\sum_{k=1}^{\infty}\frac1{k^\alpha}\) имеем:
		\[
		\sum_{k=1}^{\infty}\frac1{k^\alpha}
		\begin{cases}
			\text{сходится}, & \alpha>1,\\
			\text{расходится}, & \alpha\le 1.
		\end{cases}
		\]
		Это следует из сравнения с интегралом \(\int_1^{+\infty}x^{-\alpha}\,dx\).
	\end{enumerate}
\end{Example}

\subsection{Признак Лейбница}

\begin{Definition}{Ряд Лейбница}{def:leibniz-series}
	Ряд \(\sum_{k=1}^{\infty}a_k\) называется рядом Лейбница, если:
	\begin{enumerate}
		\item знаки соседних членов чередуются: \(\operatorname{sign}a_k=-\operatorname{sign}a_{k+1}\);
		\item \(a_k\ne 0\);
		\item \(|a_k|\) монотонно убывает;
		\item \(|a_k|\to 0\) при \(k\to\infty\).
	\end{enumerate}
\end{Definition}

\begin{Theorem}{Признак Лейбница}{thm:leibniz-test}
	Всякий ряд Лейбница сходится.
\end{Theorem}

\begin{proof}
Достаточно доказать критерий Коши. Пусть \(n,m\in\mathbb N\). Рассмотрим хвост
\[
T_{n,m}=\sum_{k=n}^{n+m}a_k.
\]
Так как знаки чередуются, а \(|a_k|\) убывает, то частичные суммы такого хвоста остаются между \(0\) и первым по модулю членом хвоста. Поэтому
\[
|T_{n,m}|\le |a_n|.
\]
Так как \(|a_n|\to 0\), то для любого \(\varepsilon>0\) найдётся \(n_0\), что при \(n\ge n_0\) имеем \(|a_n|<\varepsilon\). Следовательно, \(|T_{n,m}|<\varepsilon\) для всех \(m\). По критерию Коши ряд сходится. 
\end{proof}

\begin{Example}{Знакочередующийся гармонический ряд}{ex:alternating-harmonic}
	Ряд
	\[
	\sum_{k=1}^{\infty}\frac{(-1)^{k-1}}{k}
	\]
	сходится по признаку Лейбница, потому что \(1/k\searrow 0\). Но он не сходится абсолютно, так как \(\sum 1/k\) расходится.
\end{Example}

\subsection{Преобразование Абеля}

\begin{Lemma}{Преобразование Абеля}{lem:abel-transform}
	Пусть \(n,m\in\mathbb N\). Обозначим
	\[
	A_j=\sum_{k=n+1}^{j}a_k,\qquad j=n+1,\dots,n+m.
	\]
	Тогда
	\[
	\sum_{k=n+1}^{n+m}a_kb_k
	=
	A_{n+m}b_{n+m}
	+
	\sum_{j=n+1}^{n+m-1}A_j(b_j-b_{j+1}).
	\]
\end{Lemma}

\begin{proof}
Положим \(A_n=0\). Тогда \(a_j=A_j-A_{j-1}\). Поэтому
\[
\sum_{j=n+1}^{n+m}a_jb_j
=
\sum_{j=n+1}^{n+m}(A_j-A_{j-1})b_j.
\]
После раскрытия суммы получаем телескопическое выражение:
\[
\sum_{j=n+1}^{n+m}a_jb_j
=
A_{n+m}b_{n+m}
+
\sum_{j=n+1}^{n+m-1}A_j(b_j-b_{j+1}).
\]
Лемма доказана. 
\end{proof}

\subsection{Признаки Абеля и Дирихле}

\begin{Theorem}{Признак Абеля}{thm:abel-test}
	Пусть ряд \(\sum_{k=1}^{\infty}a_k\) сходится, а последовательность \((b_k)\) монотонна и ограничена. Тогда ряд
	\[
	\sum_{k=1}^{\infty}a_kb_k
	\]
	сходится.
\end{Theorem}



\begin{proof}
Поскольку \((b_k)\) ограничена, существует \(B>0\), такое что \(|b_k|\le B\) для всех \(k\).

Так как \(\sum a_k\) сходится, по критерию Коши для любого \(\varepsilon>0\) найдётся \(n_0\), что при \(n\ge n_0\) все хвостовые суммы \(A_j=\sum_{k=n+1}^{j}a_k\) удовлетворяют оценке \(|A_j|<\varepsilon\).

По преобразованию Абеля:
\[
\left|\sum_{k=n+1}^{n+m}a_kb_k\right|
\le
|A_{n+m}||b_{n+m}|
+
\sum_{j=n+1}^{n+m-1}|A_j||b_j-b_{j+1}|.
\]
Так как \((b_k)\) монотонна, её полное изменение на любом конечном отрезке не превосходит \(2B\). Поэтому
\[
\left|\sum_{k=n+1}^{n+m}a_kb_k\right|
\le
\varepsilon B+\varepsilon\cdot 2B
=
3B\varepsilon.
\]
Значит, хвосты ряда \(\sum a_kb_k\) стремятся к нулю равномерно по \(m\). По критерию Коши ряд сходится.
\end{proof}

\begin{Theorem}{Признак Дирихле}{thm:dirichlet-test}
	Пусть частичные суммы ряда \(\sum_{k=1}^{\infty}a_k\) ограничены, то есть существует \(C_0>0\), такое что
	\[
	\left|\sum_{k=1}^{N}a_k\right|\le C_0
	\quad \text{для всех } N.
	\]
	Пусть также \((b_k)\) монотонна и \(b_k\to 0\). Тогда ряд
	\[
	\sum_{k=1}^{\infty}a_kb_k
	\]
	сходится.
\end{Theorem}

\begin{proof}
Из ограниченности частичных сумм следует, что все конечные хвосты ряда \(\sum a_k\) ограничены: существует \(C>0\), такое что
\[
\left|\sum_{k=p}^{q}a_k\right|\le C
\quad \text{для всех } p\le q.
\]
Например, можно взять \(C=2C_0\).

Возьмём \(\varepsilon>0\). Так как \(b_k\to 0\), найдётся \(n_0\), что при \(k\ge n_0\) выполнено \(|b_k|<\varepsilon\).

Для \(n\ge n_0\) применим преобразование Абеля к хвосту:
\[
\left|\sum_{k=n+1}^{n+m}a_kb_k\right|
\le
|A_{n+m}||b_{n+m}|
+
\sum_{j=n+1}^{n+m-1}|A_j||b_j-b_{j+1}|.
\]
Здесь \(|A_j|\le C\). Кроме того, так как \((b_k)\) монотонна и стремится к нулю,
\[
\sum_{j=n+1}^{n+m-1}|b_j-b_{j+1}|
\le
|b_{n+1}|+|b_{n+m}|
<2\varepsilon.
\]
Следовательно,
\[
\left|\sum_{k=n+1}^{n+m}a_kb_k\right|
\le
C\varepsilon+2C\varepsilon
=
3C\varepsilon.
\]
По критерию Коши ряд \(\sum a_kb_k\) сходится.
\end{proof}

\begin{Example}{Применение признака Дирихле}{ex:dirichlet-example}
	Ряд \(\sum_{k=1}^{\infty}\frac{\sin k}{k}\) сходится по признаку Дирихле: частичные суммы \(\sum_{k=1}^{N}\sin k\) ограничены, а последовательность \(1/k\) монотонно убывает к нулю.
\end{Example}

\subsection{Перестановки членов ряда}

\begin{Definition}{Перестановка членов ряда}{def:series-permutation}
	Перестановкой натурального ряда называется биекция \(\sigma:\mathbb N\to\mathbb N\). Ряд
	\[
	\sum_{k=1}^{\infty} a_{\sigma(k)}
	\]
	называется перестановкой ряда \(\sum_{k=1}^{\infty}a_k\).
\end{Definition}

\begin{Theorem}{Перестановка абсолютно сходящегося ряда}{thm:absolute-permutation}
	Если ряд \(\sum_{k=1}^{\infty}a_k\) сходится абсолютно, то для любой перестановки \(\sigma:\mathbb N\to\mathbb N\) ряд
	\[
	\sum_{k=1}^{\infty}a_{\sigma(k)}
	\]
	также сходится абсолютно и имеет ту же сумму.
\end{Theorem}

\begin{proof}
Обозначим \(S=\sum_{k=1}^{\infty}a_k\). Так как \(\sum |a_k|\) сходится, то для любого \(\varepsilon>0\) найдётся \(M\), такое что
\[
\sum_{j=M+1}^{\infty}|a_j|<\varepsilon.
\]
Поскольку \(\sigma\) — биекция, среди чисел \(\sigma(1),\dots,\sigma(N)\) при достаточно большом \(N\) встретятся все номера \(1,2,\dots,M\).

Тогда при \(n\ge N\) разность между \(S\) и частичной суммой переставленного ряда состоит только из тех членов \(a_j\), у которых \(j>M\). Поэтому
\[
\left|
S-\sum_{k=1}^{n}a_{\sigma(k)}
\right|
\le
\sum_{j=M+1}^{\infty}|a_j|
<\varepsilon.
\]
Значит, \(\sum a_{\sigma(k)}\) сходится к \(S\).

Абсолютная сходимость переставленного ряда доказывается так же, применяя уже доказанное к ряду \(\sum |a_k|\). 
\end{proof}

\begin{Remark}{Почему нужна абсолютная сходимость}{rem:conditional-permutation}
	Для условно сходящихся рядов перестановки могут менять сумму ряда или даже приводить к расходимости. Поэтому в теореме о перестановках условие абсолютной сходимости существенно.
\end{Remark}

\subsection{Теорема Римана о перестановках условно сходящегося ряда}

\begin{Theorem}{Теорема Римана}{thm:riemann-rearrangement}
	Пусть ряд \(\sum_{k=1}^{\infty} a_k\) сходится условно, то есть ряд \(\sum a_k\) сходится, но ряд \(\sum |a_k|\) расходится.
	
	Тогда для любого числа \(A\in\mathbb R\) существует перестановка \(\sigma:\mathbb N\to\mathbb N\), такая что ряд
	\[
	\sum_{k=1}^{\infty} a_{\sigma(k)}
	\]
	сходится к \(A\).
\end{Theorem}

\begin{proof}
	Для простоты считаем, что \(a_k\ne 0\) для всех \(k\). Нулевые члены не влияют на сумму ряда и могут быть расставлены произвольно.
	
	Выделим отдельно положительные и отрицательные члены исходного ряда:
	\[
	a_{p_1},a_{p_2},\dots >0,
	\qquad
	a_{q_1},a_{q_2},\dots <0,
	\]
	где внутри каждой подпоследовательности порядок такой же, как в исходном ряде.
	
	Так как ряд \(\sum a_k\) сходится условно, то
	\[
	\sum_{j=1}^{\infty} a_{p_j}=+\infty,
	\qquad
	\sum_{j=1}^{\infty} a_{q_j}=-\infty.
	\]
	Иначе, если бы сумма положительных членов была конечна, то из сходимости \(\sum a_k\) следовала бы конечность суммы отрицательных членов по модулю, а значит сходился бы \(\sum |a_k|\), что противоречит условной сходимости.
	
	Теперь строим переставленный ряд. Берём положительные члены до тех пор, пока частичная сумма впервые не станет больше \(A\). Затем берём отрицательные члены до тех пор, пока частичная сумма впервые не станет меньше \(A\). Потом снова берём положительные члены до перехода выше \(A\), затем отрицательные до перехода ниже \(A\), и так далее.
	
	То есть получается ряд
	\[
	b_1+b_2+\dots,
	\qquad
	b_k=a_{\sigma(k)},
	\]
	где положительные и отрицательные члены идут блоками. Каждый шаг возможен, потому что сумма всех ещё не использованных положительных членов равна \(+\infty\), а сумма всех ещё не использованных отрицательных членов равна \(-\infty\).
	
	Докажем, что полученный ряд сходится к \(A\). Обозначим
	\[
	S_N=\sum_{k=1}^{N}b_k.
	\]
	Так как исходный ряд сходится, то \(a_k\to 0\). Поэтому и для переставленных членов тоже \(b_k\to 0\).
	
	После каждого положительного блока сумма впервые становится больше \(A\). Значит, превышение над \(A\) не больше последнего добавленного положительного члена. После каждого отрицательного блока сумма впервые становится меньше \(A\). Значит, отклонение вниз от \(A\) не больше модуля последнего добавленного отрицательного члена.
	
	Поскольку эти последние добавленные члены стремятся к нулю, для любого \(\varepsilon>0\) начиная с некоторого места все переходы через уровень \(A\) происходят с ошибкой меньше \(\varepsilon\). Внутри каждого блока частичные суммы лежат между двумя соседними переходами через \(A\), поэтому тоже отличаются от \(A\) меньше чем на \(\varepsilon\).
	
	Следовательно,
	\[
	\forall \varepsilon>0\ \exists N_0\ \forall N\ge N_0:
	|S_N-A|<\varepsilon.
	\]
	Значит, \(S_N\to A\), то есть \(\sum b_k=A\).
\end{proof}

\begin{Remark}{Смысл теоремы Римана}{rem:riemann-meaning}
	У условно сходящегося ряда сумма зависит от порядка слагаемых. Более того, перестановкой членов можно получить любую заранее заданную сумму \(A\in\mathbb R\).
	
	Это резко отличается от абсолютно сходящихся рядов: у них перестановка членов не меняет сумму.
\end{Remark}

\subsection{Линейные пространства}

\begin{Definition}{Линейное пространство над \(\mathbb R\)}{def:linear-space}
	Множество \(V\) называется линейным пространством над \(\mathbb R\), если на нём заданы операции сложения и умножения на число:
	\[
	+:V\times V\to V,
	\qquad
	\cdot:\mathbb R\times V\to V,
	\]
	и выполнены следующие свойства:
	
	\begin{enumerate}
		\item \(u+v=v+u\);
		\item \((u+v)+w=u+(v+w)\);
		\item существует нулевой вектор \(0\in V\), такой что \(v+0=v\);
		\item для каждого \(v\in V\) существует противоположный вектор \(-v\), такой что \(v+(-v)=0\);
		\item \(1\cdot v=v\);
		\item \((\lambda\mu)v=\lambda(\mu v)\);
		\item \((\lambda+\mu)v=\lambda v+\mu v\);
		\item \(\lambda(u+v)=\lambda u+\lambda v\).
	\end{enumerate}
	
	Здесь \(u,v,w\in V\), \(\lambda,\mu\in\mathbb R\).
\end{Definition}

\begin{Example}{Примеры линейных пространств}{ex:linear-spaces}
	\begin{enumerate}
		\item \(\mathbb R^n\) со стандартными операциями:
		\[
		(x_1,\dots,x_n)+(y_1,\dots,y_n)
		=
		(x_1+y_1,\dots,x_n+y_n),
		\]
		\[
		\lambda(x_1,\dots,x_n)
		=
		(\lambda x_1,\dots,\lambda x_n).
		\]
		
		\item Пространство всех многочленов с вещественными коэффициентами.
		
		\item Пространство \(C[0,1]\) всех непрерывных функций \(f:[0,1]\to\mathbb R\).
		
		\item Пространство всех вещественных последовательностей \(x=(x_1,x_2,\dots)\).
	\end{enumerate}
\end{Example}

\subsection{Нормированные пространства}

\begin{Definition}{Норма}{def:norm}
	Пусть \(V\) — линейное пространство над \(\mathbb R\). Нормой на \(V\) называется отображение
	\[
	\|\cdot\|:V\to\mathbb R_+,
	\]
	такое что для любых \(u,v\in V\) и любого \(\lambda\in\mathbb R\) выполнено:
	
	\begin{enumerate}
		\item \(\|v\|\ge 0\), причём \(\|v\|=0\iff v=0\);
		\item \(\|\lambda v\|=|\lambda|\cdot \|v\|\);
		\item \(\|u+v\|\le \|u\|+\|v\|\).
	\end{enumerate}
\end{Definition}

\begin{Definition}{Нормированное линейное пространство}{def:normed-space}
	Линейное пространство \(V\) с заданной на нём нормой \(\|\cdot\|\) называется нормированным линейным пространством.
\end{Definition}

\begin{Example}{Нормы в \(\mathbb R^n\)}{ex:norms-rn}
	Для \(x=(x_1,\dots,x_n)\in\mathbb R^n\) часто используются следующие нормы:
	\[
	\|x\|_2=\sqrt{x_1^2+\dots+x_n^2},
	\qquad
	\|x\|_{\infty}=\max_{1\le k\le n}|x_k|.
	\]
	Более общий пример:
	\[
	\|x\|_p=
	\left(\sum_{k=1}^{n}|x_k|^p\right)^{1/p},
	\qquad
	1\le p<\infty.
	\]
\end{Example}

\begin{Example}{Норма на пространстве непрерывных функций}{ex:c01-norm}
	На пространстве \(C[0,1]\) естественно задаётся равномерная норма:
	\[
	\|f\|_{\infty}
	=
	\max_{x\in[0,1]}|f(x)|.
	\]
	Максимум существует, потому что непрерывная функция на отрезке достигает своих наибольшего и наименьшего значений.
\end{Example}

\begin{Example}{Пространство многочленов}{ex:polynomials}
	Множество многочленов на \([0,1]\)
	\[
	p(x)=\sum_{k=0}^{n}a_kx^k
	\]
	является линейным пространством. На нём можно задать норму
	\[
	\|p\|_{\infty}
	=
	\max_{x\in[0,1]}|p(x)|.
	\]
	Это пространство является подпространством \(C[0,1]\).
\end{Example}

\begin{Remark}{Приближение непрерывных функций многочленами}{rem:weierstrass}
	Для любой непрерывной функции \(f\in C[0,1]\) существует последовательность многочленов \((p_n)\), такая что
	\[
	\|p_n-f\|_{\infty}
	=
	\max_{x\in[0,1]}|p_n(x)-f(x)|
	\to 0.
	\]
	Это означает, что многочлены могут равномерно приближать непрерывные функции на отрезке.
\end{Remark}

\subsection{Метрика, порождённая нормой}

\begin{Definition}{Метрика из нормы}{def:norm-metric}
	Если на \(V\) задана норма, то можно определить расстояние между векторами:
	\[
	\rho(x,y)=\|x-y\|.
	\]
	Такая метрика называется метрикой, порождённой нормой.
\end{Definition}

\begin{Statement}{Норма порождает метрику}{stmt:norm-metric}
	Функция \(\rho(x,y)=\|x-y\|\) действительно является метрикой, то есть:
	\begin{enumerate}
		\item \(\rho(x,y)\ge 0\), причём \(\rho(x,y)=0\iff x=y\);
		\item \(\rho(x,y)=\rho(y,x)\);
		\item \(\rho(x,z)\le \rho(x,y)+\rho(y,z)\).
	\end{enumerate}
\end{Statement}

\begin{proof}
	Первые два свойства следуют из свойств нормы:
	\[
	\rho(x,y)=\|x-y\|\ge 0,
	\qquad
	\|x-y\|=0\iff x-y=0\iff x=y.
	\]
	Симметрия:
	\[
	\rho(x,y)=\|x-y\|=\|-(y-x)\|=\|y-x\|=\rho(y,x).
	\]
	Неравенство треугольника:
	\[
	\rho(x,z)=\|x-z\|
	=
	\|(x-y)+(y-z)\|
	\le
	\|x-y\|+\|y-z\|
	=
	\rho(x,y)+\rho(y,z).
	\]
\end{proof}

\begin{Definition}{Сходимость в нормированном пространстве}{def:norm-convergence}
	Последовательность \((x_n)\) в нормированном пространстве \(V\) сходится к \(x\in V\), если
	\[
	\|x_n-x\|\to 0.
	\]
	Пишут \(x_n\to x\) в норме.
\end{Definition}

\subsection{Банаховы пространства}

\begin{Definition}{Полнота и банахово пространство}{def:banach}
	Нормированное пространство называется полным, если всякая фундаментальная последовательность в нём сходится к элементу этого же пространства.
	
	Полное нормированное линейное пространство называется банаховым пространством.
\end{Definition}

\begin{Example}{Банахово пространство \(C[0,1]\)}{ex:c01-banach}
	Пространство \(C[0,1]\) с нормой
	\[
	\|f\|_{\infty}=\max_{x\in[0,1]}|f(x)|
	\]
	является банаховым пространством.
	
	Смысл полноты здесь такой: если последовательность непрерывных функций равномерно фундаментальна, то она равномерно сходится к некоторой непрерывной функции.
\end{Example}

\subsection{Пространства последовательностей}

\begin{Definition}{Пространство \(\ell^p\)}{def:lp-space}
	Пусть \(1\le p<\infty\). Пространством \(\ell^p\) называется множество всех вещественных последовательностей \(x=(x_1,x_2,\dots)\), для которых
	\[
	\sum_{k=1}^{\infty}|x_k|^p<\infty.
	\]
	Норма в этом пространстве задаётся формулой
	\[
	\|x\|_p=
	\left(\sum_{k=1}^{\infty}|x_k|^p\right)^{1/p}.
	\]
\end{Definition}

\begin{Definition}{Пространство \(\ell^{\infty}\)}{def:linfty-space}
	Пространством \(\ell^{\infty}\) называется множество всех ограниченных вещественных последовательностей:
	\[
	\ell^{\infty}
	=
	\{x=(x_1,x_2,\dots):\sup_{k\in\mathbb N}|x_k|<\infty\}.
	\]
	Норма задаётся так:
	\[
	\|x\|_{\infty}
	=
	\sup_{k\in\mathbb N}|x_k|.
	\]
\end{Definition}

\begin{Remark}{Почему в \(\ell^\infty\) используется супремум}{rem:sup-linfty}
	У ограниченной последовательности максимум может не достигаться. Например, \(x_k=1-\frac1k\) ограничена сверху числом \(1\), но ни один член не равен \(1\). Поэтому в определении нормы используется \(\sup\), а не \(\max\).
\end{Remark}

\subsection{Ряды в нормированных пространствах}

\begin{Definition}{Ряд в нормированном пространстве}{def:series-normed-space}
	Пусть \(V\) — нормированное пространство, \(v_k\in V\). Рядом в \(V\) называется выражение
	\[
	\sum_{k=1}^{\infty}v_k.
	\]
	Он сходится к \(S\in V\), если частичные суммы
	\[
	S_N=\sum_{k=1}^{N}v_k
	\]
	сходятся к \(S\) по норме, то есть \(\|S_N-S\|\to 0\).
\end{Definition}

\begin{Definition}{Абсолютная сходимость в нормированном пространстве}{def:absolute-normed-space}
	Ряд \(\sum v_k\) называется абсолютно сходящимся, если сходится числовой ряд
	\[
	\sum_{k=1}^{\infty}\|v_k\|.
	\]
\end{Definition}

\begin{Theorem}{Абсолютная сходимость в банаховом пространстве}{thm:absolute-banach}
	Пусть \(V\) — банахово пространство. Если ряд \(\sum v_k\) абсолютно сходится, то он сходится в \(V\).
\end{Theorem}

\begin{proof}
	Пусть сходится числовой ряд \(\sum \|v_k\|\). Тогда по критерию Коши для числовых рядов:
	\[
	\forall \varepsilon>0\ \exists n_0\ \forall n\ge n_0\ \forall m\in\mathbb N:
	\sum_{k=n+1}^{n+m}\|v_k\|<\varepsilon.
	\]
	Для частичных сумм \(S_N=\sum_{k=1}^{N}v_k\) получаем
	\[
	\|S_{n+m}-S_n\|
	=
	\left\|\sum_{k=n+1}^{n+m}v_k\right\|
	\le
	\sum_{k=n+1}^{n+m}\|v_k\|
	<\varepsilon.
	\]
	Значит, \((S_N)\) — фундаментальная последовательность в \(V\). Так как \(V\) полно, существует \(S\in V\), такое что \(S_N\to S\). Следовательно, ряд \(\sum v_k\) сходится.
\end{proof}